% =============================================================================
%  CHAPTER 2 -- PROBABILITY, RETURNS, AND STATIONARITY FROM SCRATCH
% =============================================================================
\section{Probability, returns, and stationarity from scratch}
\label{ch:probability}

This chapter builds the entire statistical vocabulary the rest of the book uses. If you
already know it, skim; the notation is what matters, because every later derivation
refers back to it.

\subsection{Random variables, and the two kinds of average}
\label{sec:prob:rv}

A \emph{random variable} is a quantity whose value is not known in advance but whose
possible values have a description. Tomorrow's return on Adobe is a random variable;
yesterday's is a number. The description is either a \emph{probability mass function}
(discrete outcomes) or a \emph{density} $f(x)$ (continuous outcomes), with
$\int f(x)\dd x = 1$ and $f\ge 0$.

Two averages must never be confused:

\begin{description}[leftmargin=0pt,style=nextline]
  \item[The population mean (expectation)] a property of the \emph{process}:
    \begin{equation}
      \mu \;=\; \E[X] \;=\; \int x\, f(x)\dd x .
      \label{eq:expectation}
    \end{equation}
    You never observe it. It is a model parameter.
  \item[The sample mean] a property of the \emph{data} you actually collected:
    \begin{equation}
      \bar X_T \;=\; \frac{1}{T}\sum_{t=1}^{T}X_t .
      \label{eq:samplemean}
    \end{equation}
    You always observe it. It is a statistic.
\end{description}

The central activity of statistics is using the second to say something about the first,
and quantifying how much the second would wiggle if you collected different data. That
wiggle is the \emph{sampling distribution}, and it is the origin of every standard
error, confidence interval, and $p$-value in this book.

\subsection{Variance, covariance, correlation}
\label{sec:prob:variance}

\begin{definition}[Variance and standard deviation]
\label{def:variance}
$\Var(X) = \E\bigl[(X-\E X)^2\bigr] = \E[X^2]-(\E X)^2$, and
$\sigma_X=\sqrt{\Var(X)}$.
\end{definition}

The second equality is worth deriving because it is used constantly:
\[
  \E\bigl[(X-\mu)^2\bigr] = \E[X^2] - 2\mu\E[X] + \mu^2
  = \E[X^2] - 2\mu^2 + \mu^2 = \E[X^2]-\mu^2 .
\]

Variance measures dispersion in \emph{squared} units, which is why we usually quote
$\sigma$. Two algebraic facts do most of the work in this book:
\begin{align}
  \Var(aX+b) &= a^2\Var(X) , \label{eq:varaffine}\\
  \Var(X+Y) &= \Var(X)+\Var(Y)+2\Cov(X,Y) . \label{eq:varsum}
\end{align}
Equation~\eqref{eq:varsum} is the mathematical statement of ``diversification'': if
$\Cov(X,Y)<0$, the sum has less variance than the sum of variances. A hedged
stat-arb position is a sum with a deliberately negative covariance term.

\begin{definition}[Covariance and correlation]
$\Cov(X,Y)=\E[(X-\E X)(Y-\E Y)]=\E[XY]-\E X\,\E Y$, and
\begin{equation}
  \rho_{XY} \;=\; \Corr(X,Y) \;=\; \frac{\Cov(X,Y)}{\sigma_X\sigma_Y}\in[-1,1].
  \label{eq:correlation}
\end{equation}
\end{definition}

The bound $|\rho|\le1$ is the Cauchy--Schwarz inequality and it is what makes
correlation interpretable as a unit-free ``co-movement''. Correlation is invariant to
rescaling ($\Corr(aX+b,Y)=\sign(a)\rho$) but \emph{not} to nonlinear transformation, and
$\rho=0$ means only ``no linear relationship'' --- two variables can be perfectly
dependent and uncorrelated (take $X\sim\mathcal N(0,1)$, $Y=X^2$; then $\E[XY]=\E[X^3]=0$).

\paragraph{Sample versions.} Given $T$ paired observations,
\begin{equation}
  \hat\sigma_X^2 = \frac{1}{T-1}\sum_{t=1}^{T}(X_t-\bar X)^2,
  \qquad
  \widehat{\Cov}(X,Y) = \frac{1}{T-1}\sum_{t=1}^{T}(X_t-\bar X)(Y_t-\bar Y).
  \label{eq:samplevar}
\end{equation}

Why $T-1$ and not $T$? Because $\bar X$ is itself estimated from the same data, so the
deviations $(X_t-\bar X)$ are not free: they satisfy the constraint
$\sum_t (X_t-\bar X)=0$, leaving only $T-1$ independent pieces of information
(\emph{degrees of freedom}). Formally, with $s^2_T = \frac{1}{T}\sum (X_t-\bar X)^2$,
\[
  \E[s_T^2] = \E\Bigl[\frac{1}{T}\sum X_t^2 - \bar X^2\Bigr]
  = \frac{1}{T}\Bigl[T(\sigma^2+\mu^2)\Bigr] - \Bigl(\sigma^2/T + \mu^2\Bigr)
  = \sigma^2\Bigl(1-\frac{1}{T}\Bigr),
\]
so the $\frac{1}{T}$ version \emph{under}-estimates $\sigma^2$ by a factor
$(1-1/T)$; multiplying by $T/(T-1)$ removes the bias. The identical accounting
reappears in regression, where the divisor becomes $T-n$ for $n$ estimated
coefficients (Section~\ref{sec:ols:sigma2}).

\paragraph{Vectors and matrices of second moments.} For a random vector
$x\in\R^{n}$, the covariance matrix is $\Sigma=\E[(x-\E x)(x-\E x)^{\top}]\in\R^{n\times n}$:
symmetric, with $\Sigma_{ii}=\Var(x_i)$ and $\Sigma_{ij}=\Cov(x_i,x_j)$.
$\Sigma$ is \emph{positive semi-definite} (PSD): for any $a\in\R^n$,
$a^{\top}\Sigma a = \Var(a^{\top}x)\ge 0$. PSD-ness is why covariance matrices have
real non-negative eigenvalues and can be Cholesky-factorised --- the computational
backbone of Chapter~\ref{ch:ols} and of the C++ engine.

\subsection{The distributions that show up, and why}
\label{sec:prob:distributions}

\begin{description}[leftmargin=0pt,style=nextline]
  \item[Normal] $X\sim\mathcal N(\mu,\sigma^2)$ has density
    $f(x)=\frac{1}{\sigma\sqrt{2\pi}}\exp\!\bigl(-\frac{(x-\mu)^2}{2\sigma^2}\bigr)$.
    It matters because of the central limit theorem, not because returns are normal
    (they are not: they have fat tails). Sums of many small independent effects look
    Gaussian; a daily return is a sum of thousands of individual trades.
  \item[$\chi^2_k$] The sum of squares of $k$ independent standard normals,
    $\sum_{i=1}^k Z_i^2\sim\chi^2_k$, with $\E=k$, $\Var=2k$. It is the distribution of
    a \emph{variance} statistic, which is why heteroskedasticity and portmanteau tests
    (Chapter~\ref{ch:timeseries}) have $\chi^2$ limits.
  \item[$t_\nu$] $T=Z/\sqrt{V/\nu}$ with $Z\sim\mathcal N(0,1)$, $V\sim\chi^2_\nu$
    independent. It is the distribution of ``estimate divided by its own estimated
    standard error'', i.e.\ of essentially every $t$-statistic in this book. Heavier
    tails than the normal: with $\nu=30$ the 5\% two-sided critical value is 2.04, not
    1.96.
  \item[$F_{k,\nu}$] The ratio of two scaled independent $\chi^2$ variables. It is the
    distribution of ``does a group of $k$ restrictions hurt the fit?'' --- used by the
    joint-significance test and by Ramsey's RESET test
    (Section~\ref{sec:nonlinear:reset}).
\end{description}

\subsection{Prices, returns, and the adjustment arithmetic}
\label{sec:prob:returns}

Prices are the raw data; \emph{returns} are the economically meaningful quantity,
because a \$200 stock is not ``bigger'' than a \$20 stock in any way that matters for
risk.

\begin{definition}[Simple and log returns]
\begin{equation}
  R_t \;=\; \frac{P_t}{P_{t-1}}-1,
  \qquad
  r_t \;=\; \log\frac{P_t}{P_{t-1}} = \log(1+R_t).
  \label{eq:returns}
\end{equation}
\end{definition}

Log returns are used throughout this project for three concrete reasons.

\paragraph{(1) Time additivity.} Over $h$ bars,
\begin{equation}
  \log\frac{P_{t+h}}{P_t} \;=\; \sum_{k=1}^{h}\log\frac{P_{t+k}}{P_{t+k-1}}
  \;=\; \sum_{k=1}^{h} r_{t+k},
  \label{eq:logadditive}
\end{equation}
so multi-period log returns are \emph{sums}, and sums are what probability theory is
good at. Simple returns are \emph{products}, $\prod(1+R_k)-1$, which is much less
tractable.

\paragraph{(2) Symmetry.} A $+10\%$ move followed by $-10\%$ leaves you at
$1.1\times0.9=0.99$, i.e.\ $-1\%$: simple returns are not reversible. Log returns are:
$\log 1.1 + \log 0.9 = 0$. A spread built from log prices is therefore a quantity that
returns to zero when relative prices return to their starting configuration --- exactly
the property a mean-reversion signal needs.

\paragraph{(3) Local equivalence.} By the Taylor expansion of $\log(1+x)$ around 0,
\begin{equation}
  r_t \;=\; R_t - \tfrac{1}{2}R_t^2 + \tfrac{1}{3}R_t^3 - \cdots
  \;\approx\; R_t \quad \text{for } \lvert R_t\rvert \ll 1,
  \label{eq:logtaylor}
\end{equation}
and for a daily equity move ($\lvert R_t\rvert\sim 1\%$) the discrepancy is
$\sim 0.005\%$ --- negligible relative to the 10\,bps trading cost. So nothing is lost
by working in logs at daily frequency.

\begin{remark}[Honesty: levels or returns?]
Stat-arb spreads are usually built from \emph{log price levels}, not from returns,
because the claim being tested is about the \emph{relative level} of two prices
(``Adobe is expensive versus its peers''), and because the tradable object must itself
be mean-reverting. Regressing levels on levels is statistically dangerous
(Section~\ref{sec:prob:spurious}), which is precisely why the diagnostics of
Chapter~\ref{ch:timeseries} exist: we must establish that the \emph{residual} of a
levels regression is stationary before we are allowed to trade it.
\end{remark}

\paragraph{Corporate actions: why ``adjusted close''.} A 2-for-1 stock split halves the
price and doubles the share count; the holder's wealth is unchanged. An unadjusted
price series therefore shows a $-50\%$ ``return'' on the split date that is pure
accounting. The multiplicative adjustment is: for a split with ratio $k$ (new shares per
old share), all prices \emph{before} the event are divided by $k$; for a cash dividend
$D$ paid when the price is $P$, all prior prices are multiplied by $(1-D/P)$. Yahoo's
\code{Adj Close} applies both, so
\begin{equation}
  r_t^{\text{adj}} \;=\; \log\frac{P^{\text{adj}}_t}{P^{\text{adj}}_{t-1}}
  \label{eq:totalreturn}
\end{equation}
is a \emph{total} return (capital gain plus reinvested dividends). A split-adjusted-only
series gives \emph{price} returns and silently drops $\sim$1.5--2.5\%/year of dividend
income per name --- which is why Chapter~\ref{ch:phase1} uses the Yahoo panel as primary
and treats the split-adjusted-only NYSE/Kaggle panel as a secondary cross-check.

\subsection{Stationarity: the property that makes trading possible}
\label{sec:prob:stationarity}

\begin{definition}[Strict and weak stationarity]
A process $\{X_t\}$ is \emph{strictly stationary} if the joint distribution of
$(X_{t_1},\dots,X_{t_k})$ equals that of $(X_{t_1+h},\dots,X_{t_k+h})$ for every shift
$h$. It is \emph{weakly} (second-order) stationary if
$\E[X_t]=\mu$ is constant, $\Var(X_t)=\sigma^2<\infty$ is constant, and
$\Cov(X_t,X_{t+\ell})=\gamma(\ell)$ depends only on the lag $\ell$, not on $t$.
\end{definition}

Weak stationarity is what econometrics uses, and it is exactly what a mean-reversion
strategy needs: a constant mean to revert \emph{to}, a constant variance to
standardize \emph{by} (equation~\eqref{eq:saa:z}), and an autocovariance that depends
only on how far apart two times are.

\paragraph{The random walk is not stationary.} Let $X_t = X_{t-1}+\eta_t$ with
$\eta_t$ i.i.d.\ mean zero, variance $\sigma_\eta^2$. Iterating,
$X_t = X_0 + \sum_{s=1}^t \eta_s$, so
\begin{equation}
  \Var(X_t) \;=\; t\,\sigma_\eta^2 \;\longrightarrow\;\infty ,
  \qquad
  \Cov(X_t, X_{t+\ell}) = t\,\sigma_\eta^2 .
  \label{eq:rwvariance}
\end{equation}
Both the variance and the autocovariance depend on $t$: the process has no stable
``typical distance from the mean'', and there is no mean to revert to. Stock prices are,
to a first approximation, random walks with drift --- which is the formal version of
``you cannot predict the market''. We say such a series is integrated of order one,
$I(1)$: it becomes stationary after differencing once. A series that is already
stationary is $I(0)$.

\begin{intuition}
The whole game of statistical arbitrage, in one sentence: \emph{prices are $I(1)$, so
find a linear combination of several $I(1)$ prices that is $I(0)$, and trade that
combination.} Such a combination is called \emph{cointegrated}
(Section~\ref{sec:ts:cointegration}). The residual $\eps_t$ of
equation~\eqref{eq:basemodel} is the candidate, and Chapters~\ref{ch:timeseries}
and~\ref{ch:ou} are devoted to testing whether it really is $I(0)$ and how fast it
comes home.
\end{intuition}

\paragraph{Autocovariance and autocorrelation functions.} For a weakly stationary
series, $\gamma(\ell)=\Cov(X_t,X_{t+\ell})$ and
$\rho(\ell)=\gamma(\ell)/\gamma(0)$. The sample autocorrelation at lag $\ell$ is
\begin{equation}
  \hat\rho_\ell \;=\; \frac{\sum_{t=\ell+1}^{T}(X_t-\bar X)(X_{t-\ell}-\bar X)}
                          {\sum_{t=1}^{T}(X_t-\bar X)^2}.
  \label{eq:acf}
\end{equation}
Under the null of white noise, $\hat\rho_\ell$ is approximately
$\mathcal N(0,1/T)$, so the familiar $\pm 1.96/\sqrt T$ bands on an ACF plot are not
decoration --- they are a 5\% test per lag. The Durbin--Watson and Ljung--Box statistics
of Section~\ref{sec:ts:autocorr} are just disciplined summaries of the same object.

\subsection{Spurious regression: why levels are dangerous}
\label{sec:prob:spurious}

Take two \emph{independent} random walks $Y_t$ and $X_t$. Regress one on the other. You
will routinely see $R^2$ above 0.5 and $t$-statistics above 3, on data generated with no
relationship whatsoever. This is not a quirk, it is a theorem
\cite{grangernewbold1974}: the OLS $t$-statistic diverges as $T\to\infty$
(it grows like $\sqrt T$), so with enough data any two trending series look
``significantly related''.

The mechanism is easy to see. Both series are dominated by their accumulated trends; the
regression matches trend to trend; the residual is the difference of two random walks,
which is itself usually a random walk --- non-stationary, so its sample variance does
not settle, and the usual formulas for standard errors (which assume i.i.d.\ residuals)
are simply invalid.

Two defences, and this project uses both:
\begin{enumerate}
  \item \textbf{Test the residual for a unit root.} If the residual of a levels
    regression is $I(0)$, the regression is not spurious --- it is a cointegrating
    regression. This is the Engle--Granger procedure
    (Section~\ref{sec:ts:cointegration}).
  \item \textbf{Work in differences/returns and accept a smaller signal.} Always valid,
    but it throws away the level information that a mean-reversion trade needs.
\end{enumerate}

This is the reason the diagnostics chapter is not optional decoration: without it, the
strategy's beautiful $R^2\approx0.98$ on log prices would be indistinguishable from a
spurious regression.

\subsection{Convergence: LLN, CLT, and the origin of $\sqrt T$}
\label{sec:prob:clt}

\begin{theorem}[Weak law of large numbers]
If $X_1,\dots,X_T$ are i.i.d.\ with mean $\mu$ and finite variance, then
$\bar X_T \xrightarrow{\ p\ } \mu$.
\end{theorem}

\begin{theorem}[Central limit theorem]
Under the same conditions,
\begin{equation}
  \sqrt{T}\,(\bar X_T-\mu)\ \xrightarrow{\ d\ }\ \mathcal N(0,\sigma^2),
  \qquad\text{equivalently}\qquad
  \bar X_T \approx \mathcal N\Bigl(\mu,\ \frac{\sigma^2}{T}\Bigr).
  \label{eq:clt}
\end{equation}
\end{theorem}

Everything about statistical power follows from the $\sqrt T$ in \eqref{eq:clt}: to halve
the standard error of an average you need \emph{four times} as much data. This is why a
sixteen-year daily panel (4{,}190 bars) is not ``a lot of data'' for detecting a small
daily edge. If your daily mean P\&L is $\mu_d$ with daily standard deviation
$\sigma_d$, the $t$-statistic for ``$\mu_d>0$'' is
\begin{equation}
  t \;=\; \frac{\bar d}{\hat\sigma_d/\sqrt{T}}
  \;=\; \frac{\mu_d}{\sigma_d}\sqrt{T}
  \;=\; \frac{\text{daily Sharpe}}{1}\cdot\sqrt{T},
  \label{eq:sharpetstat}
\end{equation}
so a daily Sharpe of $0.003$ (which annualizes to $\approx0.05$) needs
$T\approx(1.96/0.003)^2\approx 427{,}000$ days --- about 1{,}700 years --- to be
significant at 5\%. \emph{This single calculation explains the entire negative result of
this project}, and it is quantified properly in Chapters~\ref{ch:backtestmath}
and~\ref{ch:multipletesting}.

For dependent data the CLT still holds but $\sigma^2$ must be replaced by the
\emph{long-run variance}
\begin{equation}
  \sigma^2_{LR} \;=\; \gamma(0) + 2\sum_{\ell=1}^{\infty}\gamma(\ell)
  \;=\; \lim_{T\to\infty} T\,\Var(\bar X_T),
  \label{eq:longrunvar}
\end{equation}
which is what HAC estimators (Section~\ref{sec:ts:hac}) estimate. If the series is
positively autocorrelated, $\sigma^2_{LR}>\gamma(0)$ and naive standard errors are
\emph{too small} --- the standard way a backtest overstates its own significance.

\subsection{Estimators: bias, variance, and the MSE decomposition}
\label{sec:prob:estimators}

An \emph{estimator} $\hat\theta$ is a function of the data intended to approximate an
unknown parameter $\theta$. Being a function of random data, it is itself random.

\begin{definition}[Bias, variance, MSE, consistency]
\[
  \operatorname{Bias}(\hat\theta)=\E[\hat\theta]-\theta,\qquad
  \Var(\hat\theta)=\E\bigl[(\hat\theta-\E\hat\theta)^2\bigr],\qquad
  \operatorname{MSE}(\hat\theta)=\E\bigl[(\hat\theta-\theta)^2\bigr].
\]
$\hat\theta$ is \emph{consistent} if $\hat\theta\xrightarrow{\ p\ }\theta$ as
$T\to\infty$, and \emph{efficient} if it attains the smallest possible variance among
unbiased estimators (the Cramér--Rao bound).
\end{definition}

\begin{proposition}[Bias--variance decomposition]
\begin{equation}
  \operatorname{MSE}(\hat\theta) \;=\; \Var(\hat\theta) + \operatorname{Bias}(\hat\theta)^2 .
  \label{eq:biasvariance}
\end{equation}
\end{proposition}
\begin{proof}
Let $b=\E\hat\theta-\theta$. Then
$\E[(\hat\theta-\theta)^2]=\E[(\hat\theta-\E\hat\theta + b)^2]
=\E[(\hat\theta-\E\hat\theta)^2] + 2b\E[\hat\theta-\E\hat\theta] + b^2
=\Var(\hat\theta)+b^2$, since the middle term is zero.
\end{proof}

\begin{intuition}
Equation~\eqref{eq:biasvariance} is the licence to be biased. An unbiased estimator can
have enormous variance (OLS with 100 correlated regressors and 120 observations); a
slightly biased estimator with much smaller variance can have a lower MSE. Ridge
regression (Chapter~\ref{ch:regularization}) is exactly this trade made deliberately:
accept bias, buy variance, and the tuning parameter $\lambda$ is the dial between them.
The same decomposition governs the nonlinear-vs-linear verdict of
Chapter~\ref{ch:nonlinear}: gradient boosting has almost no bias and enormous variance
at $T\approx3{,}000$ daily bars, and loses.
\end{intuition}

\subsection{Hypothesis testing, from nothing}
\label{sec:prob:testing}

A test has five ingredients:

\begin{enumerate}
  \item A \textbf{null hypothesis} $H_0$ (the boring, default state of the world) and an
    \textbf{alternative} $H_1$.
  \item A \textbf{test statistic} $\hat\tau$, a number computed from the data whose
    distribution \emph{under $H_0$} is known or can be simulated.
  \item A \textbf{significance level} $\alpha$ (conventionally 0.05).
  \item A \textbf{decision rule}: reject $H_0$ if $\hat\tau$ falls in the critical
    region $C_\alpha$, chosen so that $P(\hat\tau\in C_\alpha\mid H_0)=\alpha$.
  \item A \textbf{$p$-value}: $p=\sup_{\theta\in H_0}P(\hat\tau \text{ at least as
    extreme as observed}\mid\theta)$, i.e.\ the smallest $\alpha$ at which you would
    reject.
\end{enumerate}

\begin{remark}[What a $p$-value is not]
$p$ is \emph{not} the probability that $H_0$ is true, and $1-p$ is not the probability
that your strategy works. It is the probability of seeing data this extreme \emph{if
nothing were going on}. The two errors are: Type~I (reject a true null; probability
$\alpha$, the ``false discovery'') and Type~II (fail to reject a false null;
probability $\beta$, the ``missed discovery''). \emph{Power} is $1-\beta$, and it grows
with effect size, with $T$, and with $\alpha$. Testing 44 strategies at $\alpha=0.05$
guarantees roughly $44\times0.05\approx2.2$ false discoveries among the nulls --- which
is why Chapter~\ref{ch:multipletesting} exists.
\end{remark}

\paragraph{Worked example: the one-sample $t$-test.} Data $X_1,\dots,X_T$ i.i.d.\
$\mathcal N(\mu,\sigma^2)$; test $H_0:\mu=0$ vs.\ $H_1:\mu\ne0$. The statistic is
\[
  \hat\tau \;=\; \frac{\bar X}{s/\sqrt T},\qquad
  s^2=\frac{1}{T-1}\sum(X_t-\bar X)^2 .
\]
Under $H_0$: $\bar X\sim\mathcal N(0,\sigma^2/T)$ so $\sqrt T\bar X/\sigma\sim\mathcal
N(0,1)$; and $(T-1)s^2/\sigma^2\sim\chi^2_{T-1}$ independently of $\bar X$ (a special
property of the normal family). Their ratio is therefore exactly
$t_{T-1}$-distributed. Reject if $|\hat\tau|>t_{T-1,1-\alpha/2}$. If $\sigma$ were
known, the statistic would be standard normal --- the entire difference between the
normal and $t$ distributions is the price of estimating $\sigma$ from the same data.

\subsection{Annualization: the $\sqrt{252}$ rule and when it lies}
\label{sec:prob:annualization}

There are $\approx252$ U.S.\ trading days per year. If daily returns were i.i.d.\ with
mean $\mu_d$ and variance $\sigma_d^2$, then by \eqref{eq:logadditive} the annual return
is a sum of 252 of them, so
\begin{equation}
  \mu_a = 252\,\mu_d,
  \qquad
  \sigma_a^2 = 252\,\sigma_d^2,
  \qquad
  \mathrm{SR}_a = \frac{\mu_a}{\sigma_a} = \sqrt{252}\,\frac{\mu_d}{\sigma_d}
  = \sqrt{252}\,\mathrm{SR}_d .
  \label{eq:annualization}
\end{equation}
That is the origin of the ``multiply the Sharpe ratio by $\sqrt{252}$'' rule.

It lies when returns are autocorrelated, and the correction is worth deriving because
it is used in Chapter~\ref{ch:multipletesting}. With autocovariances $\gamma(\ell)$,
\[
  \Var\Bigl(\sum_{t=1}^{252} r_t\Bigr)
  = 252\,\gamma(0) + 2\sum_{\ell=1}^{251}(252-\ell)\,\gamma(\ell)
  = 252\,\gamma(0)\Bigl[1 + 2\sum_{\ell=1}^{251}\Bigl(1-\frac{\ell}{252}\Bigr)\rho(\ell)\Bigr],
\]
so
\begin{equation}
  \mathrm{SR}_a \;=\; \sqrt{252}\,\mathrm{SR}_d\,\Bigl[\,1+2\textstyle\sum_{\ell\ge1}
  \bigl(1-\tfrac{\ell}{252}\bigr)\rho(\ell)\Bigr]^{-1/2}.
  \label{eq:loadjust}
\end{equation}
This is \cite{lo2002statistics}'s adjustment. Positive autocorrelation (momentum)
inflates the naive annualized Sharpe; negative autocorrelation (mean reversion ---
our case) \emph{deflates} it, so the naive $\sqrt{252}$ figure is conservative here.
Since our holding period is $\approx$16--18 bars, $\rho(\ell)$ is materially non-zero
out to lag $\sim20$, and this book reports both the naive and the caveat.

\subsection{Linear algebra you will need in the next chapter}
\label{sec:prob:linalg}

\paragraph{Vectors, inner products, norms.} For $a,b\in\R^n$:
$a^{\top}b=\sum_i a_ib_i$ (a scalar; the inner product),
$\lVert a\rVert_2=\sqrt{a^{\top}a}$, $\lVert a\rVert_1=\sum_i\lvert a_i\rvert$,
$\lVert a\rVert_\infty=\max_i\lvert a_i\rvert$. The three norms define the three
geometries of Chapter~\ref{ch:regularization}: L2 (ridge), L1 (lasso), L$\infty$.

\paragraph{Matrix calculus.} The two identities used over and over:
\begin{equation}
  \nabla_{\beta}\,(a^{\top}\beta) = a,
  \qquad
  \nabla_{\beta}\,(\beta^{\top}A\beta) = (A+A^{\top})\beta
  \;\overset{A=A^{\top}}{=}\; 2A\beta .
  \label{eq:matrixcalculus}
\end{equation}
Also $\nabla_\beta\lVert Y-X\beta\rVert_2^2 = -2X^{\top}(Y-X\beta)$, obtained by
expanding $Y^{\top}Y-2\beta^{\top}X^{\top}Y+\beta^{\top}X^{\top}X\beta$ and applying
\eqref{eq:matrixcalculus} with $A=X^{\top}X$ (symmetric).

\paragraph{Eigen-decomposition and the spectral theorem.} If $A\in\R^{n\times n}$ is
symmetric, there is an orthogonal $V$ ($V^{\top}V=I$) and a diagonal
$\Lambda=\diag(\lambda_1,\dots,\lambda_n)$ with
\begin{equation}
  A \;=\; V\Lambda V^{\top} \;=\; \sum_{i=1}^n \lambda_i\, v_i v_i^{\top},
  \qquad A v_i = \lambda_i v_i .
  \label{eq:spectral}
\end{equation}
For a PSD matrix all $\lambda_i\ge0$. This decomposition is simultaneously: the
definition of principal components (Chapter~\ref{ch:regularization}), the way ridge
shrinkage is understood (it damps small $\lambda_i$ directions), the reason the Kalman
covariance stays PSD, and the reason a cyclic-Jacobi eigen-solver exists in the C++
engine (Section~\ref{sec:arch:principles}).

\paragraph{Condition number.} $\kappa(A)=\lambda_{\max}/\lambda_{\min}$ for SPD $A$.
It bounds how much the solution of $A\beta=b$ can move relative to a perturbation of
$b$: $\lVert\delta\beta\rVert/\lVert\beta\rVert \le \kappa(A)\,\lVert\delta b\rVert/\lVert b\rVert$.
Correlated regressors make $\kappa(X^{\top}X)$ huge, so the estimated weights swing
wildly on tiny data changes --- the numerical face of multicollinearity, and the
motivation for the ridge penalty $\lambda I$ that raises every eigenvalue by $\lambda$.

\paragraph{Positive definiteness and Cholesky.} $A$ SPD means $a^{\top}Aa>0$ for all
$a\ne0$; then $A=LL^{\top}$ with $L$ lower-triangular and positive diagonal, computed in
$n^3/3$ flops without pivoting (Section~\ref{sec:ols:cholesky}). Solving
$A\beta=b$ becomes two triangular solves, $Lz=b$ then $L^{\top}\beta=z$, each $O(n^2)$,
and is roughly twice as fast and twice as accurate as forming $A^{-1}$ explicitly.
\emph{We never form the inverse} --- not in the math, and not in the C++.
